\section{Data, Variables, and Methods}\label{sec:data}

\subsection{Study Area and Data}

The analysis covers 30 provincial-level units in mainland China from 2013 to 2021. Tibet, Hong Kong, Macao, and Taiwan are excluded because of data availability and comparability. The panel is balanced with 270 observations. The sample starts in 2013 because the National Bureau of Statistics introduced a unified urban-rural household survey that year, which makes the income series comparable over the period.

Provincial economic, social, infrastructure, and environmental indicators are taken from the \emph{China Statistical Yearbook} (2014--2022), the \emph{China City Statistical Yearbook} (2014--2022), the \emph{China Environmental Statistics Yearbook}, and provincial yearbooks \citep{nbs2022,nbscity2022}. Digital financial inclusion sub-indices come from the Peking University Digital Financial Inclusion Index \citep{pku2021}. A small number of missing values were filled by linear interpolation within provinces. For urban and rural incomes, we checked the extracted values against the accounting identity linking total disposable income to urban and rural incomes and the urbanization rate, and corrected cells that violated it.

\subsection{Digital Factor Allocation Index}

DFA is measured with a four-dimension system (Table~\ref{tab:indicators}). Digital infrastructure captures the network foundations that determine whether digital services reach firms and households; digital industry captures the scale and market application of digital production; digital finance captures the coverage, intensity, and digitalization of financial services; and digital human capital captures the supply of workers who can use digital technologies.

The 15 indicators are selected for relevance to the dimension, availability across all provinces and years, and consistency of definitions. To prevent a few large provinces from dominating the index, indicators are expressed as per capita, per GDP, or per area ratios. DFA assigns equal weights to the four dimensions and entropy weights to indicators within each dimension, so each factor category has the same ex ante importance while the data determine the weights inside it.

\subsection{Regional Coordinated Development Index}

RCD is measured with a three-dimension system (Table~\ref{tab:indicators}). Economic coordination uses GDP per capita, retail sales per capita, and fiscal self-sufficiency. Social coordination uses disposable income per capita, hospital beds per 10,000 people, higher-education students per 100,000 people, and the share of social security spending. Urban-rural coordination uses the urbanization rate, the urban-rural income ratio (negative), and the urban registered unemployment rate (negative).

The three dimensions receive equal weights and entropy weights are applied within each dimension. For robustness we also construct a five-dimension index that adds infrastructure and ecological coordination, a coupling coordination degree for the three core dimensions, and a global entropy-weighted index.

\begin{table}[htbp]
\caption{Indicator systems for digital factor allocation (DFA) and regional coordinated development (RCD).}\label{tab:indicators}
\footnotesize
\begin{tabular}{>{\raggedright\arraybackslash}p{3.0cm}>{\raggedright\arraybackslash}p{4.2cm}>{\raggedright\arraybackslash}p{5.0cm}c}
\toprule
Dimension & Indicator & Definition and unit & Sign \\
\midrule
\multicolumn{4}{l}{\textit{Panel A: Digital factor allocation (DFA)}} \\
\midrule
Digital infrastructure & Internet access ports per capita & Internet broadband access ports / year-end population (ports per person) & + \\
 & Optical cable density & Long-distance optical cable length / land area (km/km$^2$) & + \\
 & Mobile phone penetration & Mobile phone subscribers per 100 people & + \\
 & IPv4 addresses per capita & IPv4 addresses / year-end population & + \\
 & Domain names per capita & Registered domain names / year-end population & + \\
Digital industry & Software revenue / GDP & Software business revenue / GDP (\%) & + \\
 & E-commerce sales / GDP & E-commerce transaction value / GDP (\%) & + \\
 & Firms with e-commerce & Share of enterprises with e-commerce transactions (\%) & + \\
 & Websites per 100 firms & Number of websites per 100 enterprises & + \\
 & Information firms per capita & Information enterprises / year-end population & + \\
Digital finance & DFI coverage breadth & Peking University Digital Financial Inclusion Index sub-index & + \\
 & DFI usage depth & Peking University Digital Financial Inclusion Index sub-index & + \\
 & DFI digitization & Peking University Digital Financial Inclusion Index sub-index & + \\
Digital human capital & ICT employment share & Employment in information transmission, computer services, and software / total urban employment (\%) & + \\
 & Internet penetration & Internet users / year-end population (\%) & + \\
\midrule
\multicolumn{4}{l}{\textit{Panel B: Regional coordinated development (RCD)}} \\
\midrule
Economic coordination & GDP per capita & Regional GDP / year-end population (CNY/person) & + \\
 & Retail sales per capita & Retail sales of consumer goods / year-end population (CNY/person) & + \\
 & Fiscal self-sufficiency & Local fiscal revenue / local fiscal expenditure (\%) & + \\
Social coordination & Disposable income per capita & All residents' per capita disposable income (CNY/person) & + \\
 & Hospital beds & Hospital beds per 10,000 people & + \\
 & Higher-education students & Higher-education students per 100,000 people & + \\
 & Social security spending & Social security and employment expenditure / fiscal expenditure (\%) & + \\
Urban-rural coordination & Urbanization rate & Urban resident population / year-end population (\%) & + \\
 & Urban-rural income ratio & Urban disposable income / rural disposable income & $-$ \\
 & Urban unemployment rate & Registered urban unemployment rate (\%) & $-$ \\
\bottomrule
\end{tabular}

\noindent\footnotesize Notes: All indicators are compiled from the \emph{China Statistical Yearbook}, the \emph{China City Statistical Yearbook}, provincial statistical yearbooks, and the Peking University Digital Financial Inclusion Index. DFA assigns equal weights to the four dimensions and entropy weights to indicators within each dimension. RCD assigns equal weights to the three dimensions and entropy weights to indicators within each dimension.
\end{table}

\subsection{Entropy Weight Method}

The entropy method assigns larger weights to indicators with greater dispersion. For indicator $j$ and observation $i$, the raw value $x_{ij}$ is winsorized at the 1st and 99th percentiles and normalized. For positive indicators,
\begin{equation}
z_{ij}=\frac{x_{ij}-\min_i x_{ij}}{\max_i x_{ij}-\min_i x_{ij}},
\end{equation}
and for negative indicators,
\begin{equation}
z_{ij}=\frac{\max_i x_{ij}-x_{ij}}{\max_i x_{ij}-\min_i x_{ij}}.
\end{equation}
A constant ($10^{-4}$) is added to avoid zeros. The normalized values become proportions $p_{ij}=z_{ij}/\sum_i z_{ij}$, and the entropy of indicator $j$ is
\begin{equation}
e_j=-\frac{1}{\ln N}\sum_{i=1}^{N}p_{ij}\ln p_{ij},
\end{equation}
where $N$ is the number of observations. The weight is
\begin{equation}
w_j=\frac{1-e_j}{\sum_{k}(1-e_k)},
\end{equation}
and the composite score is $S_i=\sum_j w_j z_{ij}$. Entropy weights are computed within each dimension and the dimension scores are averaged with equal weights, which prevents a dimension with many indicators from receiving disproportionate weight.

\subsection{Baseline and Threshold Models}

The baseline specification is a two-way fixed-effects model:
\begin{equation}
RCD_{it}=\alpha_0+\alpha_1 DFA_{it}+\gamma'X_{it}+\mu_i+\delta_t+\varepsilon_{it},
\label{eq:baseline}
\end{equation}
where $i$ indexes provinces, $t$ indexes years, $\mu_i$ and $\delta_t$ are province and year fixed effects, $X_{it}$ is a vector of controls, and $\varepsilon_{it}$ is clustered at the province level.

To test H2 and H3 we estimate a Hansen (1999) panel threshold model:
\begin{equation}
RCD_{it}=\beta_0+\beta_1 DFA_{it}\cdot I(q_{it}\leq\gamma)
+\beta_2 DFA_{it}\cdot I(q_{it}>\gamma)
+\gamma'X_{it}+\mu_i+\varepsilon_{it},
\label{eq:threshold}
\end{equation}
where $q_{it}$ is the threshold variable and $I(\cdot)$ is the indicator function. The threshold is estimated by grid search over the inner 80\% of the distribution of $q$. The null of linearity is tested with a cluster wild bootstrap that generates $y^*=\hat{y}+\varepsilon^*$, where $\varepsilon^*$ is the fitted residual multiplied by a province-level Rademacher sign, and re-estimates the threshold and the $F$ statistic in each replication. The $p$-value is the share of bootstrap statistics above the observed one. We also report the likelihood-ratio confidence interval for $\gamma$ using the critical value 7.35 \citep{hansen1999,hansen2000}.

The threshold variable is economic development (log GDP per capita) for H2 and financial development for H3. We also estimate specifications using DFA, technological investment, and private economy vitality as threshold variables to check whether the nonlinearity is specific to economic development.

\subsection{Controls and Descriptive Statistics}

The controls capture local conditions that may affect RCD but are not part of the index. Government intervention is local expenditure over GDP; foreign direct investment is log realized FDI per capita; population density is log year-end population per square kilometer; and private economy vitality is private enterprises per 10,000 people. Together they cover the fiscal, external, spatial, and market dimensions of local conditions.

Table~\ref{tab:descriptives} reports descriptive statistics. DFA has a mean of 0.259 and a standard deviation of 0.179; RCD has a mean of 0.357 and a standard deviation of 0.148. The urban-rural income ratio averages 2.53 and ranges from 1.84 to 3.56.

\begin{table}[htbp]
\caption{Descriptive statistics.}\label{tab:descriptives}
\begin{threeparttable}
\small
\begin{tabular}{lrrrrrr}
\toprule
Variable & N & Mean & SD & Min & Median & Max \\
\midrule
DFA index & 270 & 0.259 & 0.179 & 0.007 & 0.223 & 1.000 \\
RCD index & 270 & 0.357 & 0.148 & 0.083 & 0.333 & 0.846 \\
RCD5 index & 270 & 0.414 & 0.117 & 0.137 & 0.406 & 0.805 \\
RCD coupling coordination & 270 & 0.555 & 0.180 & 0.051 & 0.550 & 1.000 \\
Urban-rural income ratio & 270 & 2.533 & 0.361 & 1.842 & 2.474 & 3.556 \\
Government intervention & 270 & 26.292 & 11.122 & 10.501 & 23.207 & 75.344 \\
Log FDI per capita & 270 & 18.762 & 1.250 & 15.540 & 18.402 & 21.650 \\
Log population density & 270 & 5.469 & 1.289 & 2.068 & 5.659 & 8.282 \\
Private economy vitality & 270 & 191.904 & 142.191 & 41.889 & 155.839 & 813.382 \\
Financial development & 270 & 7.604 & 3.072 & 3.241 & 6.892 & 19.634 \\
Technological investment & 270 & 2.204 & 1.511 & 0.539 & 1.559 & 6.757 \\
Industrial structure upgrading & 270 & 140.959 & 74.470 & 66.533 & 122.565 & 524.401 \\
\bottomrule
\end{tabular}
\begin{tablenotes}
\small
\item Notes: The sample covers 30 provinces from 2013 to 2021 (N = 270). DFA is the digital factor allocation index; RCD is the regional coordinated development index; RCD5 is the five-dimension alternative; RCD-CCD is the coupling coordination degree of the three core dimensions. See Table~\ref{tab:indicators} for definitions.
\end{tablenotes}
\end{threeparttable}
\end{table}

\subsection{Spatial Weights and Spatial Specifications}\label{sec:spatial}

The spatial analysis uses the 30 provinces as units. The weight matrix $W$ is a row-standardized rook contiguity matrix built from the provincial boundary file. Hainan has no land neighbour and is linked to Guangdong, the conventional treatment. As a robustness check we also use an inverse-squared-distance matrix from provincial centroids. Both have a zero diagonal, so $W$ times a variable is the weighted average of the other provinces in the same year. Spatial autocorrelation is measured with the global Moran's $I$ and with local indicators (LISA), each with 999 random permutations; a province is classified as a local cluster when the pseudo $p$-value is below 0.10.

We estimate three spatial extensions of the threshold model. The first adds the spatial lag of the regressor,
\begin{equation}
\begin{aligned}
RCD_{it}={}&\beta_1 DFA_{it}\cdot I(q_{it}\leq\gamma)
+\beta_2 DFA_{it}\cdot I(q_{it}>\gamma)\\
&+\delta\,(W\!\cdot\!DFA)_{it}+\gamma'X_{it}+\mu_i+\varepsilon_{it},
\end{aligned}
\label{eq:slx}
\end{equation}
which is not identified here: after the within transformation DFA and its spatial lag are correlated at 0.979, so the two coefficients cannot be separated with 30 provinces. We report it as a transparency check only.

The second is a spatial error model,
\begin{equation}
\begin{aligned}
RCD_{it}={}&\beta_1 DFA_{it}\cdot I(q_{it}\leq\gamma)
+\beta_2 DFA_{it}\cdot I(q_{it}>\gamma)
+\gamma'X_{it}+\mu_i+u_{it},\\
&u_{it}=\lambda\sum_j w_{ij}u_{jt}+e_{it},
\end{aligned}
\label{eq:sem}
\end{equation}
and the third is a spatial autoregressive model,
\begin{equation}
\begin{aligned}
RCD_{it}={}&\rho\sum_j w_{ij}RCD_{jt}
+\beta_1 DFA_{it}\cdot I(q_{it}\leq\gamma)\\
&+\beta_2 DFA_{it}\cdot I(q_{it}>\gamma)
+\gamma'X_{it}+\mu_i+\varepsilon_{it}.
\end{aligned}
\label{eq:sar}
\end{equation}
Both are estimated by maximum likelihood with province fixed effects. For a given spatial parameter the concentrated likelihood is maximized over $\gamma$ by grid search, with the log-determinant evaluated from the eigenvalues of $W$. The threshold test is a cluster wild bootstrap conditional on the estimated spatial parameter.
